\section{Topological Manifolds}

\noindent\textbf{Exercise~\ref{ex:topologicalmfds-1}.}\label{sol:topologicalmfds-1}

The M\"{o}bius strip can be represented by the following diagram
\begin{center}
    \btik
        \draw[thick] (0,0) -- (5,0);
        \draw[thick] (0,2) -- (5,2);
        \draw[thick, decoration={markings, mark=at position 0.5 with {\arrow{>}}}, postaction={decorate}] (0,0) -- (0,2);
        \draw[thick, decoration={markings, mark=at position 0.5 with {\arrow{>}}}, postaction={decorate}] (5,2) -- (5,0);
        \draw[blue, ultra thick] (0,1.5) .. controls (2,0) and (4,1.5) ..(5,0.5);
    \etik
\end{center}
where the arrows indicate how we identify the two edges together (i.e. the bottom right corner goes to the top left corner). The blue line is the river. It is clear that we will need at least two charts in order to map the whole strip. We could choose them as follows

\begin{center}
    \btik
        \draw[dashed, fill=lightgray] (0,0) -- (0,2) -- (1.5,2) -- (1.5,0) -- (0,0);
        \draw[dashed, fill=lightgray] (5,0) -- (5,2) -- (3.5,2) -- (3.5,0) -- (5,0);
        \draw[thick] (0,0) -- (5,0);
        \draw[thick] (0,2) -- (5,2);
        \draw[thick, decoration={markings, mark=at position 0.5 with {\arrow{>}}}, postaction={decorate}] (0,0) -- (0,2);
        \draw[thick, decoration={markings, mark=at position 0.5 with {\arrow{>}}}, postaction={decorate}] (5,2) -- (5,0);
        \draw[blue, ultra thick] (0,1.5) .. controls (2,0) and (4,1.5) ..(5,0.5);
        \node at (0.75,0.5) {\Large{$U_1$}};
        \node at (4.25,1.5) {\Large{$U_1$}};
        %
        \draw[dashed, fill=lightgray] (7.5,0) -- (7.5,2) -- (11.5,2) -- (11.5,0) -- (7.5,0);
        \draw[thick] (7,0) -- (12,0);
        \draw[thick] (7,2) -- (12,2);
        \draw[thick, decoration={markings, mark=at position 0.5 with {\arrow{>}}}, postaction={decorate}] (7,0) -- (7,2);
        \draw[thick, decoration={markings, mark=at position 0.5 with {\arrow{>}}}, postaction={decorate}] (12,2) -- (12,0);
        \draw[blue, ultra thick] (7,1.5) .. controls (9,0) and (11,1.5) ..(12,0.5);
        \node at (9.5,1.5) {\Large{$U_2$}};
        %
    \etik
\end{center}
We then define maps from these shaded regions into $\R^2$. For example we can just map them so that they are the shaded regions on the page, taking care to region the two parts of $U_1$ properly. To save my writing more Tikz code, I'll leave this to your imagination.

\bigskip
\noindent\textbf{Exercise~\ref{ex:topologicalmfds-2}.}\label{sol:topologicalmfds-2}

The map is $x: F^1 \to [-1,1]$ given by $x:(m,n)
\mapsto m$. Clearly this map is not injective as $x\big((m,n)\big) = x\big((m,n')\big)$ where $n' = -n$.

\bigskip
\noindent\textbf{Exercise~\ref{ex:topologicalmfds-3}.}\label{sol:topologicalmfds-3}

The problem was that $n'=-n$ gave non-injectivity, so clearly we define
\bse
    \begin{split}
        U_{\uparrow}  & := \{(m,n)\in\R^2 \, | \, m^4 + n^4 = 1, n>0\} \\
        U_{\downarrow}  & := \{(m,n)\in\R^2 \, | \, m^4 + n^4 = 1, n<0\}
    \end{split}
\ese
and so have the maps
\bse
    x_{\uparrow} : U_{\uparrow} \to (-1,1), \qquad \text{and} \qquad x_{\downarrow} : U_{\downarrow} \to (-1,1).
\ese
It is important that we take the inequalities for $n$ (i.e. $n\neq 0$) so that our sets are open, as required by the question.

\bigskip
\noindent\textbf{Exercise~\ref{ex:topologicalmfds-4}.}\label{sol:topologicalmfds-4}

Same as above we simply define
\bse
    V_{\uparrow}  := \{(m,n)\in\R^2 \, | \, m^4 + n^4 = 1, m>0\},
\ese
and the map
\bse
    y_{\uparrow} : V_{\uparrow} \to (-1,1),
\ese
given by $y_{\uparrow}:(m,n)\mapsto n$.

\bigskip
\noindent\textbf{Exercise~\ref{ex:topologicalmfds-5}.}\label{sol:topologicalmfds-5}

Yes it is invertible as it is bijective. We construct the inverse as
\bse
    y_{\uparrow}^{-1} : (-1,1) \to V_{\uparrow},
\ese
where the action is given by
\bse
    y_{\uparrow}^{-1} : n \mapsto \big( \sqrt[4]{1-n^4}, n\big).
\ese

\bigskip
\noindent\textbf{Exercise~\ref{ex:topologicalmfds-6}.}\label{sol:topologicalmfds-6}

Yes their domains overlap as
\bse
    U_{\uparrow}\cap V_{\uparrow} = \{(m,n)\in\R^2 \, | \, m^4 + n^4 = 1, n,m>0\}
\ese
The transition map is
\bse
    \big(x_{\uparrow}\circ y_{\uparrow}^{-1}\big) : (0,1) \to (0,1)
\ese
with
\bse
    \big(x_{\uparrow}\circ y_{\uparrow}^{-1}\big)(a) = x_{\uparrow}\big(\sqrt[4]{1-a^4}, a\big) = \sqrt[4]{1-a^4} \in (0,1).
\ese

\bigskip
\noindent\textbf{Exercise~\ref{ex:topologicalmfds-7}.}\label{sol:topologicalmfds-7}

The answer is obviously 4 maps, $x_{\uparrow},x_{\downarrow},y_{\uparrow}$ and $y_{\downarrow}$, and yes they form an atlas. From this we see that $F^1$ is a topological manifold of dimension $d=1$.
